%% USPSC-Resumo.tex
\setlength{\absparsep}{18pt} % ajusta o espaçamento dos parágrafos do resumo		
\begin{resumo}
	\begin{flushleft} 
			\setlength{\absparsep}{0pt} % ajusta o espaçamento da referência	
			\SingleSpacing 
			\imprimirautorabr~~\textbf{\imprimirtituloresumo}.	\imprimirdata. \pageref{LastPage} p. 
			%Substitua p. por f. quando utilizar oneside em \documentclass
			%\pageref{LastPage} f.
			\imprimirtipotrabalho~-~\imprimirinstituicao, \imprimirlocal, \imprimirdata. 
 	\end{flushleft}
\OnehalfSpacing 			
% >>> ESCREVA AQUI O SEU RESUMO <<<
% O resumo deve ressaltar objetivo, método, resultados e conclusões (ABNT NBR 6028).
% Texto único, em parágrafo único, entre 150 e 500 palavras.
Escreva aqui o resumo do seu trabalho.

 \textbf{Palavras-chave}: palavra-chave 1. palavra-chave 2. palavra-chave 3.
\end{resumo}
